\documentclass[10pt,a4paper]{amsart}
\usepackage[margin=19mm]{geometry}
\usepackage{amsmath,amssymb,mathtools}
\usepackage[scr=rsfs,cal=euler]{mathalfa}
\usepackage{xfrac}
\usepackage[unicode,hidelinks,bookmarks=false]{hyperref}
\newcommand{\ie}{i.e.\ }
\newcommand{\cf}{cf.\ }
\newcommand{\resp}{respectively\ }
\newcommand{\Subsection}{\subsection}
\providecommand{\nobreakdash}{\nobreak\hbox{-}}
\setlength{\emergencystretch}{2em}
\multlinegap0pt
\usepackage{amssymb}
\usepackage{enumerate}
\usepackage{algorithm}\usepackage{algpseudocode}
\usepackage{dsfont}
\usepackage[shortlabels]{enumitem}
\usepackage{multirow}
\usepackage{booktabs}
\usepackage{subfig}
\usepackage{xcolor}
\newtheorem{proposition}{Proposition}
\usepackage{cleveref}
\crefname{proposition}{Proposition}{Propositions}
\newcommand{\ares}[2]{\operatorname{ares}_{#2}^{#1}}
\newcommand{\bres}[2]{\operatorname{bres}_{#2}^{#1}}
\DeclareMathOperator{\id}{Id}
\title{Anisotropic Proximal Point Algorithm}
\author{Emanuel Laude, Panagiotis Patrinos}
\date{Source: arXiv:2312.09834v4 (CC BY 4.0)}
\begin{document}
\maketitle

\noindent\emph{Source and licence.} Anisotropic Proximal Point Algorithm. Version arXiv:2312.09834v4 (CC BY 4.0), distributed by arXiv under the Creative Commons Attribution 4.0 International licence, http://creativecommons.org/licenses/by/4.0/, as recorded in the arXiv licence metadata for this exact version and shown on the abstract page https://arxiv.org/abs/2312.09834v4. Full licence notice: \texttt{licenses/arxiv-2312.09834-CC-BY-4.0.txt}.\par\smallskip

\noindent\textit{Editorial scope.} Counted unit: the article's proposition thm:absorb\_relaxation\_stepsize (source0.tex lines 612--618). The article's complete original proof of it is delivered with it (source0.tex lines 619--636, one \texttt{proof} environment of 18 lines). No mathematics is written, repaired or re-derived: the statement and the proof are the article's own lines, in the article's own notation, and no retained line is substituted or re-flowed.  Retained uncounted as the source-local context the proof needs: lines 539--545.  Nothing was omitted from any retained span, so the package carries no omission record.  No retained line contains a citation, so the wrapper carries no bibliography and no bibliography entry.  No retained line contains a figure environment, an image-inclusion command or drawing code, so no image is delivered and the package declares no asset dependency.  Editorial formatting only, in addition: an AMS \texttt{amsart} preamble that restates the article's own declarations and macros used by the retained lines, namely the environments \texttt{proposition} on separate counters, together with the standard packages those lines need.  The retained environments print in this document as Proposition 1 in order of appearance, whereas the article prints the same items with the numbering of its own full document.  The complete native source of the article as distributed in the arXiv e-print for this exact version is bundled unchanged as \texttt{sources/151556/source0.tex}.
\begin{proposition}[Moreau's decomposition] \label{thm:moreau_decomposition}
Let $T: X \rightrightarrows X^*$ be a set-valued mapping. Then the following identity holds:
\begin{align}
\ares{\phi}{T}  = \id - \nabla \phi^* \circ \bres{\phi^*}{T^{-1}} \circ \nabla \phi.
\end{align}
If, furthermore, $T$ is monotone, $\ares{\phi}{T}$ is single-valued.
\end{proposition}

\begin{proposition} \label{thm:absorb_relaxation_stepsize}
    Let $T: X \rightrightarrows X^*$ be a set-valued mapping and let $\rho \geq 0$ and $\tau >0$. Then the following identity holds:
    \begin{equation}
    (\id + \tau \nabla \phi^* \circ T_\rho)^{-1} = (1 - \lambda) \id + \lambda(\id + \gamma\nabla \phi^* \circ T)^{-1},
    \end{equation}
    for $\gamma = \tau + \rho$ and $\lambda = \tfrac{\tau}{\gamma} \in (0, 1]$.
\end{proposition}

\begin{proof}
Note that $T_\gamma = \bres{\gamma\phi^*}{T^{-1}} \circ \nabla (\gamma \star\phi)$.
Invoking \cref{thm:moreau_decomposition} and reminding ourselves that $(\gamma \star \phi)^* = \gamma \phi^*$ we have:
\begin{align} \label{eq:identity_tgamma}
\gamma \nabla \phi^*(T_\gamma(x))&= (x - (\id + \gamma \nabla \phi^* \circ T)^{-1}(x)).
\end{align}
Define $T_\rho^{-1}:= (T_\rho)^{-1} = \rho \nabla \phi^* + T^{-1}$.
Then we have for any $x\in X$ the identity
\begin{align*}
(1 - \lambda) x + \lambda(\id + \gamma \nabla \phi^* \circ T)^{-1}(x) &= x - \lambda(x - (\id + \gamma \nabla \phi^* \circ T)^{-1}(x)) \\
&= x - \tau \nabla \phi^*(T_\gamma(x)) \\
&= x - \tau \nabla \phi^*\big((\tau \nabla \phi^* + (\rho \nabla \phi^* + T^{-1}))^{-1}(x)\big) \\
&= x - \tau \nabla \phi^*\big((\tau \nabla \phi^* + T_\rho^{-1})^{-1}(x)\big) \\
&= x - \tau \nabla \phi^*\big(\bres{\tau \phi^*}{T_\rho^{-1}}(\nabla (\tau \star \phi)(x))\big) \\
&=(\id + \tau \nabla \phi^* \circ T_\rho)^{-1}(x),
\end{align*}
where the second identity follows from \cref{eq:identity_tgamma} since $\lambda \gamma = \tau $, the third by $\gamma = \tau+\rho$ and the definition of $T_\gamma$, the fourth by definition of $T_\rho^{-1}$ and the last by \cref{thm:moreau_decomposition}.
\end{proof}

\end{document}
